\section{Conclusions}
In this work, I have explored the use of MSE, cosine similarity, PCA informed LDA, PCA informed XGBoost, standard XGB and ensemble SCNN methods for the classification of colon cancer samples as non-cancerous, hyperplastic, adenomatous, or cancerous via the analysis of FTIR spectral data. The SCNN method was found to be the most effective, achieving an OWA of $58.7\pm3.7\%$ when performing four-class classification of on a four-fold cross-validated test set. When mapped to binary classification, the SCNN also performs best and achieves non-cancer vs cancer classification with OWA $83.2\pm5.9\%$ - similar to the accuracy achieved by Lasalvia et al. \cite{lasalvia_discrimination_2023}. This model was then explained using LIME and CRIME techniques to extract multiple contexts to explain each tissue classification. The spectra and weightings associated with these contexts were also used to create LCRIME models. The best of these LCRIME models for 4-class classification utilised a $20\%$ supervised CP-VAE. This model was capable of reproducing SCNN predictions with OWA $47.5\pm7.6\%$ on the four-fold cross-validated test set. For binary classification, the best LCRIME model utilised a $20\%$ supervised DT-VAE and achieved an OWA of $74.7\pm 10.9\%$ on the same test set.\\ 

\noindent
Future works could improve the performance of the classification models by implementing ordinal regressive losses \cite{} during SCNN training, implementing transformer networks directly for classification, or exploring the use of different weak learner ensemble methods instead of simple voting \cite{}. Furthermore, the explainability work could be expanded on by utilising different dataset folds to achieve finer hyperparameter tuning for the CRIME VAEs, and by exploring the use of other explainability methods like SHAP \cite{}.